\setlength{\parindent}{0pt}
% ------------------------------------------------------------
\chapter{Design of Coppula}
\label{chap:design}
% ------------------------------------------------------------
\section{Pipeline Overview}
\label{sec:pipeline-overview}

\toolname\ takes an arithmetic circuit over $\mathbb{F}_2$, where the wires
are partitioned into secret, random, and public inputs
(Chapter~\ref{chap:preliminaries}), together with a probing order $d$, and
returns a verdict: the circuit is $d$-probing secure, or it is not and a
witnessing probe is handed back. Underneath sits the single question of
Section~\ref{sec:d-probing}, is every set of at most $d$ wires
secret-independent, and two difficulties have to be resolved for an answer.
Deciding one single probe is already non-trivial --- it is a statement
about a distribution over the randomness space --- and there are combinatorially many
probes to decide. The design addresses both difficulties.

Three components divide the labour. \emph{Component~I}, the rewrite engine
(Section~\ref{sec:rewrite-engine}), is the oracle for a single probe. It decides
secret-independence not by summing over a distribution but by rewriting the
probe's expression, one randomness substitution at a time, into a form with no secret in
sight (Section~\ref{sec:randomness-substitution}). \emph{Component~II}, the
subset-space traversal (Section~\ref{sec:subset-space}), is the driver. It sweeps
every $d$-subset of the wires without enumerating them one by one, in the
recursive set-splitting style of \texttt{maskVerif}
\parencite{10.1007/978-3-030-29959-0_15}. \emph{Component~III}, the probe-set
extension (Section~\ref{sec:extension}), is the amplifier. It turns the
certificate for one probe into a certificate for a strictly larger safe set of
wires, so that a single discharge settles many probes at once.

All three components someway use couplings. When Component~II reaches a batch of
probes it cannot dismiss as a whole, it asks Component~I to certify a
representative; a successful discharge returns not just a secure verdict but a
measure-preserving relabelling of the randomness space (a coupling function) that
witnesses it. Component~III substitutes that coupling in the circuit's other
wires and gathers those that are blinded (rid of secrets) into the safe set. Component~II then
splits the batch about that safe set and recurses, with the safe part already
certified. Soundness of the whole rests on Component~I alone
(Theorem~\ref{thm:rewrite-soundness}); Components~II and III only utilize its
verdicts, adding neither unsoundness nor any incompleteness of their own.

The rest of the chapter takes the three components in turn, asking
what each computes and why it is correct. The concrete
data structures that make them time- and memory-wise efficient are explained in
Chapter~\ref{chap:implementation}.

\section{Component I: The Rewrite Engine}
\label{sec:rewrite-engine}

Given a probe $\mathbf{w} = (w_1, \dots, w_k)$, $k \le d$ --- a tuple of
wires of the circuit --- the rewrite engine decides whether $\mathbf{w}$ is
secret-independent. It does so constructively, rather than reasoning about
probability distributions directly; it searches for a finite sequence of
applications of Lemma~\ref{lem:substitution}
(Section~\ref{sec:randomness-substitution}) after which $\mathbf{w}$'s form is
syntactically free of secrets --- no secret is present in the expression of
any $w_i$. Termination with a secret-free image certifies $\mathbf{w}$.
Theorem~\ref{thm:rewrite-soundness} in Section~\ref{sec:coupling-view} makes precise
why this certificate is sound.

The engine has three tiers of rewriting rules, increasing in generality and cost,
summarized in
Algorithm~\ref{alg:rewrite-engine} and mirrored in the structure of this section.
Its single entry point is the routine \textproc{CheckProbe} of
Algorithm~\ref{alg:rewrite-engine}, which tries the tiers in turn and returns a
verdict together with the derivation $T$ that Component~III later uses.
\textproc{SimpleRewrite} (Section~\ref{sec:simple-rule}) handles the common case of
every random occurring at most once, cheaply. It is tried first on the unfactored
probe and, on failure, once more after the multiplicative factoring of
Section~\ref{sec:factoring} has had a chance to expose additional single-occurrence
randoms. This factoring \& rewriting is the second tier.
\textproc{GeneralRewrite} (Section~\ref{sec:general-substitution}) is the
complete fallback for the remaining case, for random variables occurring more than once. It
substitutes a random by its context wherever the random occurs, interleaving
factoring lazily whenever both of its own find-rules stall. Section~\ref{sec:coupling-view}
closes the section by reading the whole derivation as an explicit, replayable
bijection on the randomness space. Which becomes the \gls{coupling} that Component III
(Section~\ref{sec:extension}) later utilizes in every other wire in the circuit.

\begin{algorithm}
\caption{Discharging a probe (Component I)}
\label{alg:rewrite-engine}
\small
\begin{minipage}{\linewidth}
\begin{algorithmic}[1]
\Function{CheckProbe}{$\mathbf{w}$}
  \Require probe $\mathbf{w}$, a tuple of wires of the circuit
  \Ensure a verdict, \Secure{} or \Insecure{}, with the recorded derivation $T$ (the shears of \S\ref{sec:coupling-view})
  \State $(\mathit{ok}, T) \gets \Call{SimpleRewrite}{\mathbf{w}, \langle\rangle}$ \Comment{Tier 1: unfactored}
  \If{$\mathit{ok}$} \Return $(\Secure{}, T)$ \EndIf
  \State $\mathbf{w}_f \gets \Call{Factor}{\mathbf{w}}$ \Comment{Tier 2: factor (\S\ref{sec:factoring})}
  \State $(\mathit{ok}, T) \gets \Call{SimpleRewrite}{\mathbf{w}_f, \langle\rangle}$ \Comment{retry the simple rule}
  \If{$\mathit{ok}$} \Return $(\Secure{}, T)$ \EndIf
  \State $(\mathit{ok}, T) \gets \Call{GeneralRewrite}{\mathbf{w}, \emptyset, \mathit{fuel}, \mathit{false}, \langle\rangle}$ \Comment{Tier 3: complete fallback}
  \State \Return $\mathit{ok}\ ?\ (\Secure{}, T) : (\Insecure{}, T)$
\EndFunction
\Statex
\Function{SimpleRewrite}{$\mathbf{w}, T$}
  \State $\mathit{todo} \gets \{\, r \in \mathrm{vars}(\mathbf{w}) : r \text{ occurs once, additively}, r \notin \mathbf{w} \,\}$
  \State $\mathit{todo} \gets \Call{SortByMulDepth}{\mathit{todo}}$ \Comment{shallowest multiplicative depth first}
  \While{$\mathit{todo} \neq \langle\rangle$ \textbf{and not} \Call{SecretFree}{$\mathbf{w}$}\footnote{\textproc{SecretFree}$(\mathbf{w})$ returns true exactly when no secret variable
  occurs in the current expression of any wire of~$\mathbf{w}$.}}
    \State $r \gets$ pop shallowest element of $\mathit{todo}$
    \State $x \gets$ unique additive parent of $r$ \Comment{$x = e \oplus r$}
    \State mark $x$ as a fresh random leaf
    \State $T \gets T.\mathrm{push}(r, x)$
    \State drop $r$ \Comment{may expose a new single-occurrence random for $\mathit{todo}$}
  \EndWhile
  \State \Return $(\Call{SecretFree}{\mathbf{w}}, T)$
\EndFunction
\Statex
\Function{GeneralRewrite}{$\mathbf{w}, \mathit{used}, \mathit{fuel}, \mathit{factored}, T$}
  \If{\Call{SecretFree}{$\mathbf{w}$}} \Return $(\mathit{true}, T)$ \EndIf
  \If{$\mathit{fuel} = 0$} \Return $(\mathit{false}, T)$ \EndIf
  \State $r \gets \Call{FindSimple}{\mathbf{w}}$ \Comment{single-occurrence, non-root, additive}
  \If{$r = \bot$} \ $r \gets \Call{FindGeneral}{\mathbf{w}, \mathit{used}}$ \Comment{multi-occurrence, at least one additive instance} \EndIf
  \If{$r = \bot$}
    \If{$\mathit{factored}$} \Return $(\mathit{false}, T)$ \EndIf
    \State $\mathbf{w}_f \gets \Call{Factor}{\mathbf{w}}$ \Comment{lazy factoring (\S\ref{sec:factoring})}
    \State \Return $\Call{GeneralRewrite}{\mathbf{w}_f, \mathit{used}, \mathit{fuel}, \mathit{true}, T}$
  \EndIf
  \State $x \gets$ chosen additive parent of $r$; \ $e \gets \bigoplus(\mathrm{children}(x) \setminus \{r\})$ \Comment{$r \notin \mathrm{vars}(e)$}
  \State $t \gets e \oplus r$; \ $\mathbf{w}' \gets \mathbf{w}[r \mapsto t]$ \Comment{Lemma~\ref{lem:substitution}}
  \State \Return $\Call{GeneralRewrite}{\mathbf{w}', \mathit{used} \cup \{r\}, \mathit{fuel}-1, \mathit{false}, T.\mathrm{push}(r,t)}$
\EndFunction
\end{algorithmic}
\end{minipage}
\end{algorithm}

  \subsection{The Simple Rewrite Rule}
  \label{sec:simple-rule}

  The cheapest tier discharges the pattern
  Corollary~\ref{cor:single-occurrence-discharge} addresses directly, where a random
  occurs exactly once below the probe, additively (as an operand of an XOR gate).
  Such a random is a simple-rule candidate, provided it is not itself a probe
  root. Writing $x = e \oplus r$ for its unique parent, the corollary lets us relabel
  $x$ as a fresh random and drop $r$ without ever looking at the context $e$. The
  pattern certifies that $x$, jointly with everything independent of $r$, is
  distributed as a fresh, independent randomness. The tier applies this rewrite
  repeatedly, and with each rewrite it possibly creates further randoms fitting the pattern,
  until either no secret variable is reachable from the probe (success) or no candidate
  remains.

  The exclusion of probe roots is a soundness condition. A
  root's value is what the adversary observes, so it is fixed, and relabelling
  it would claim independence for something the probe actually exposes.

  The order in which rewrites fire is based on the multiplicative depth of the random variable.
  Randoms closest to a root get rewritten first.

% refer to the soundness proof
  This tier is sound, since each step is a direct instance of
  Corollary~\ref{cor:single-occurrence-discharge}, but it is incomplete. It has nothing to
  offer for a random occurring twice, such as a mask reused across two cross-terms of a
  masked multiplication. Section~\ref{sec:general-substitution}'s fallback handles
  such cases.

  \subsection{Fallback: General Substitution}
  \label{sec:general-substitution}

  \textproc{GeneralRewrite} handles randoms that occur more than once by performing
  the substitution of Lemma~\ref{lem:substitution} literally, rather than only its
  single-occurrence case (Corollary~\ref{cor:single-occurrence-discharge}). It picks
  a random $r$ and an additive occurrence $x_1 = e \oplus r$ with
  $r \notin \mathrm{vars}(e)$, then rewrites every root by substituting
  $r \mapsto e \oplus r$ throughout. Cancellation in $\mathbb{F}_2$ collapses the
  chosen occurrence $x_1$ to the bare leaf $r$, exactly as in
  Corollary~\ref{cor:single-occurrence-discharge}, while every other occurrence of
  $r$ absorbs $e$ instead of vanishing, cancelling further if it happens to share
  structure with $e$. The loop succeeds once no secret variable remains in
  the rewritten roots.

  Each round re-examines the current probe for an applicable substitution.
  \textproc{FindSimple} looks first for the same single-occurrence, non-root pattern
  as Section~\ref{sec:simple-rule}, since a general substitution can turn a
  multi-occurrence random into a single-occurrence one; it prefers the shallowest
  candidate. Failing that, \textproc{FindGeneral} looks for any random, not yet used
  by a previous general step, occurring at least twice with some additive parent
  $x_1$ whose other children do not themselves depend on it --- the side condition
  $r \notin \mathrm{vars}(e)$ that Lemma~\ref{lem:substitution} requires.

  When both find-rules stall on the probe's current form, the engine factors
  (Section~\ref{sec:factoring}) once and retries, rather than declaring failure
  immediately. Factoring can move a random from a purely multiplicative position,
  where neither find-rule can see it, into an additive one, but on its own it
  eliminates no secret, so it must be interleaved with substitution rather than run
  once up front or once at the end. Factoring is symmetrically not tried before the
  simple rule either. Algorithm~\ref{alg:rewrite-engine} always tries
  \textproc{SimpleRewrite} unfactored first, because factoring can destroy a
  single-occurrence opportunity the unfactored form exposes, so trying the cheap,
  unfactored tier first is not just faster on the common case but strictly more
  complete.

  \textproc{FindGeneral} is needed precisely when a random occurs additively in two
  places that do not telescope into one another under the simple rule. Suppose $r$
  occurs as an operand of $x_1 = e_1 \oplus r$ in one root and as $x_2 = e_2 \oplus r$
  in another, with $r \notin \mathrm{vars}(e_1) \cup \mathrm{vars}(e_2)$.
  Substituting through $x_1$, i.e.\ $r \mapsto e_1 \oplus r$, turns $x_1$ into
  $e_1 \oplus (e_1 \oplus r) = r$ by cancellation, as in
  Corollary~\ref{cor:single-occurrence-discharge}, while the same substitution turns
  $x_2$ into $e_2 \oplus e_1 \oplus r$. If $e_1 \oplus e_2$ is itself secret-free,
  both roots are secret-free after this one step; otherwise the loop continues,
  substituting further or factoring to expose a new opportunity. This is the shape of
  the case a masked multiplication typically produces, where a single random masks
  two different cross-terms: neither term's random is disposable on its own, but the
  substitution above can still cancel the secret-dependent part they share. Whether a
  given instance is discharged this way, and in how many rounds, depends entirely on
  how much of $e_1 \oplus e_2$ the rest of the derivation can still cancel; the loop
  offers no guarantee beyond trying. Even granting unlimited fuel, it only composes
  single-coordinate shears $\widehat{\sigma}$, so it can only certify
  secret-independence witnessed by a chain of substitutions, possibly interleaved
  with factoring. Also, it should be noted that the completeness of this approach
  is proved only in the linear/affine case.
  % TODO: the proof for the linear/affine case is not given.

  Termination is only argued informally, not proved. Each general substitution
  retires its random into $\mathit{used}$ and is never repeated for it, bounding the
  number of general steps by $|\mathcal{R}|$; between general steps, simple
  substitutions strictly reduce the candidate pool; factoring is retried at most once
  per stall before the engine gives up on the current form. Interleaving factoring
  with two mutually triggering find-rules resists a clean potential-function
  argument, so a generous fuel counter backstops the loop instead. Exhausting fuel is
  treated as failure, which is sound (Theorem~\ref{thm:rewrite-soundness} is never
  invoked on a fuel-out) but not necessarily tight: a fuel-out reports the probe as
  though it were genuinely insecure, without exhibiting an actual distinguishing
  observation.

  \subsection{Relation to the Coupling-Construction View}
  \label{sec:coupling-view}

  A discharge does more than answer \Secure{}/\Insecure{}. Read the
  right way, it hands back the object Chapter~\ref{chap:couplings} asks for: a
  coupling for the probe. We state that object declaratively first, then show how
  the engine produces one.

  We let the secrets range, rather than staying fixed as in
  Section~\ref{sec:randomness-substitution}. Write
  $\mathbf{w}(\rho; \mathbf{a})$ for the value of the probe under randomness
  assignment $\rho \in \{0,1\}^{\mathcal{R}}$ and secret assignment
  $\mathbf{a} \in \{0,1\}^{\mathcal{S}}$, the public inputs are held fixed.

  \begin{definition}[Coupling of a probe]
  \label{def:probe-coupling}
  A \emph{coupling} of a probe $\mathbf{w}$ is a family
  $\{\,\Phi_{\mathbf{a},\mathbf{b}}\,\}_{\mathbf{a},\mathbf{b} \in \{0,1\}^{\mathcal{S}}}$
  of measure-preserving bijections of $\{0,1\}^{\mathcal{R}}$, one per ordered
  pair of secret assignments, satisfying
  \[
    \mathbf{w}(\rho; \mathbf{a}) \;=\; \mathbf{w}\big(\Phi_{\mathbf{a},\mathbf{b}}(\rho);\, \mathbf{b}\big)
    \qquad \text{for every } \rho \in \{0,1\}^{\mathcal{R}} .
  \]
  \end{definition}

  Definition~\ref{def:probe-coupling} says precisely where to send the randomness
  so that a run under $\mathbf{a}$ is reproduced, observation for observation, by
  a run under $\mathbf{b}$. Each $\Phi_{\mathbf{a},\mathbf{b}}$ dictates the
  relabelling of the random variables that absorb the change of secrets. Since it is
  measure-preserving, pushing the uniform randoms through it turns this pointwise
  identity into an equality of laws, so the distribution of
  $\mathbf{w}(\cdot;\mathbf{a})$ matches that of $\mathbf{w}(\cdot;\mathbf{b})$. A
  coupling existing for \emph{every} pair $(\mathbf{a},\mathbf{b})$ is thus the
  same statement as $\mathbf{w}$ being secret-independent --- which is what makes
  it the right certificate. The family carries the structure of a groupoid:
  one may take $\Phi_{\mathbf{a},\mathbf{a}} = \mathrm{id}$,
  $\Phi_{\mathbf{b},\mathbf{a}} = \Phi_{\mathbf{a},\mathbf{b}}^{-1}$, and
  $\Phi_{\mathbf{b},\mathbf{c}} \circ \Phi_{\mathbf{a},\mathbf{b}} =
  \Phi_{\mathbf{a},\mathbf{c}}$, so for any secret assignment $\mathbf{a}_0$,
  $\{\Phi_{\mathbf{a}_0,\,\cdot}\}$ already generate the rest of the couplings. It is these
  properties we care about, not any one procedure that realises them.

  What the engine records is actually such a coupling. Discharging a probe amounts to
  finding a \emph{derivation} of a finite sequence of rewrites
  \[
    \sigma_1 = (r_1 \mapsto e_1 \oplus r_1), \ \dots, \ \sigma_K = (r_K \mapsto e_K \oplus r_K),
  \]
  where $r_i \notin \mathrm{vars}(e_i)$. And after the rewrites no secret variable
  remain. It does not matter which rule proposed it, and whether the random it retires
  occurred once or many times, each $\sigma_i$ is a substitution in the sense of
  Definition~\ref{def:substitution}, hence also a measure-preserving shear
  $\widehat{\sigma_i}$ (Lemma~\ref{lem:substitution}). The coupling is assembled
  from these shears and from nothing else, so the construction below never needs
  to know how any individual rewrite was found.

  \begin{theorem}[Soundness of discharge]
  \label{thm:rewrite-soundness}
  If \textup{\textproc{CheckProbe}}$(\mathbf{w})$
  (Algorithm~\ref{alg:rewrite-engine}) returns \textup{\Secure{}}, then the
  joint distribution of $\mathbf{w}(\rho; \mathbf{a})$, for $\rho$ sampled
  uniformly at random, does not depend on the secret assignment $\mathbf{a}$.
  Moreover the recorded substitutions realise a coupling of $\mathbf{w}$ in the
  sense of Definition~\ref{def:probe-coupling}.
  \end{theorem}

  \begin{proof}
  Fix a secret assignment $\mathbf{a}$. Each recorded rewrite step contributes the
  shear $\widehat{\sigma_i}$, whose defining context $e_i$ may itself contain
  secrets, so write $\widehat{\sigma_i}^{\,\mathbf{a}}$ for its action under
  $\mathbf{a}$ and set
  \[
    \varphi_{\mathbf{a}} \;=\; \widehat{\sigma_1}^{\,\mathbf{a}} \circ
    \widehat{\sigma_2}^{\,\mathbf{a}} \circ \cdots \circ \widehat{\sigma_K}^{\,\mathbf{a}},
  \]
  which is a measure-preserving bijection of $\{0,1\}^{\mathcal{R}}$. It is the ``rewrite recipe''
  instantiated at $\mathbf{a}$.

  Also let $\overline{\mathbf{w}} = \mathbf{w}\sigma_1 \cdots
  \sigma_K$ be the probe state the derivation ends on. Applying
  Lemma~\ref{lem:substitution} once per step and composing gives, for every $\rho$,
  \[
    \mathbf{w}\big(\varphi_{\mathbf{a}}(\rho);\, \mathbf{a}\big) \;=\; \overline{\mathbf{w}}(\rho) .
  \]

  A \Secure{} verdict means $\overline{\mathbf{w}}$ is syntactically
  secret-free, so its value is
  one and the same for every secret assignment. Since
  $\varphi_{\mathbf{a}}$ preserves the uniform law on $\{0,1\}^{\mathcal{R}}$, the
  normalisation identity makes the law of $\rho \mapsto \mathbf{w}(\rho;
  \mathbf{a})$ equal to the law of $\overline{\mathbf{w}}$, which does not depend on
  $\mathbf{a}$. Hence the joint distribution of the probe does not depend on the
  secret assignment.

  For the coupling consider any two secret assignments
  $\mathbf{a}, \mathbf{b}$. The normalisation identity sends both to the
  common normal form,
  \[
    \mathbf{w}\big(\varphi_{\mathbf{a}}(\rho); \mathbf{a}\big)
    \;=\; \overline{\mathbf{w}}(\rho) \;=\;
    \mathbf{w}\big(\varphi_{\mathbf{b}}(\rho); \mathbf{b}\big) .
  \]
  Substituting $\rho \mapsto \varphi_{\mathbf{a}}^{-1}(\rho)$ and setting
  $\Phi_{\mathbf{a},\mathbf{b}} = \varphi_{\mathbf{b}} \circ
  \varphi_{\mathbf{a}}^{-1}$ yields $\mathbf{w}(\rho; \mathbf{a}) =
  \mathbf{w}\big(\Phi_{\mathbf{a},\mathbf{b}}(\rho); \mathbf{b}\big)$ for all
  $\rho$. Each $\Phi_{\mathbf{a},\mathbf{b}}$ is a composition of
  measure-preserving bijections, so the family $\{\Phi_{\mathbf{a},\mathbf{b}}\}$
  is the sought after coupling of $\mathbf{w}$.
  \end{proof}

  % I commented the following two paragraphs as it might be confusing at this point.
  % The composite $\varphi_{\mathbf{a}}$ is displayed above as a product of shears
  % only for exhibition; it is never used. Component III (Section~\ref{sec:coupling-extension})
  % uses a single member of the family as a \emph{blinding certificate}: it takes any
  % other wire $u$ in the circuit and admits it into the certified \emph{safe set}
  % when $u \circ \varphi_{\mathbf{a}}$ is secret-free --- when the same relabelling
  % that blinds the probe blinds $u$ too.
  % Evaluating $\varphi_{\mathbf{a}}$ on a sampled tape is a back-substitution
  % through the recorded steps in reverse order $i = K, \dots, 1$, since the net
  % image of $r_i$ is $t_i = e_i \oplus r_i$ evaluated against the assignment
  % already updated by every later step: the derivation performs forward
  % elimination, and reading the net map off is back-substitution through the same
  % steps in reverse.

  % Keeping this record costs nothing. Each rewrite already names its target
  % $t_i = e_i \oplus r_i$ --- the very wire it substitutes for $r_i$ --- so the
  % derivation $\sigma_1, \dots, \sigma_K$ can be accumulated as the rewrites happen,
  % and the coupling is read off that one list. Component III replays these shears
  % directly against every other wire in the circuit; no probe and no rewrite is
  % special-cased, and each step that helped certify the probe is a step that can help
  % blind its neighbours.

\section{Component II: Traversing the Subset Space}
\label{sec:subset-space}

$d$-probing security asks that \emph{every} set of at most $d$ wires be
secret-independent (Section~\ref{sec:d-probing}). Read literally, that is
$\sum_{i=1}^{d} \binom{N}{i}$ separate questions for a circuit of $N$ wires.
Component II is the
bookkeeping that keeps the question exact while never enumerating a subset
redundantly.

Two elementary observations are helpful to understand the traversal.
First, we should see that secret-independence passes to
subsets, specifically a subset of a secret-independent set is a marginal of a
secret-independent law, hence itself is secret-independent. Certifying a set
of wires therefore certifies all of its subsets in one stroke. 
Second, it is enough to look at probes of size
\emph{exactly} $d$: a smaller probe is a marginal of some size-$d$ probe
(whenever $N \ge d$, which always holds here), so it is covered once the size-$d$
probes are. The task is thus to certify every $d$-element subset of the wire set,
in batches as large as possible rather than one at a time.

  \subsection{The Representation of the Subset Space}
  \label{sec:represent-subsets}

  The traversal never holds an explicit list of subsets; it holds a compact
  description of a whole family of them at once. A \emph{factor} is a pair
  $(c, W)$ of a count $c \in \mathbb{N}$ and a wire set $W$, standing for the
  family $\binom{W}{c}$ of all $c$-element subsets of $W$. A \emph{worklist} is a
  finite sequence of factors
  \[
    wl = \big\langle (c_1, W_1), \dots, (c_m, W_m) \big\rangle,
    \qquad W_1, \dots, W_m \text{ pairwise disjoint,}
  \]
  denoting the probes obtained by drawing $c_j$ wires from each $W_j$ and taking
  the union:
  \[
    \llbracket wl \rrbracket \;=\;
      \Big\{\, S_1 \cup \dots \cup S_m \;:\; S_j \in \tbinom{W_j}{c_j} \,\Big\},
    \qquad
    \big|\llbracket wl \rrbracket\big| \;=\; \prod_{j=1}^{m} \binom{|W_j|}{c_j}.
  \]
  Disjointness of the $W_j$ makes that union a disjoint one, so every probe has
  size $\sum_j c_j$ and none is described twice. Two vacaous worklists items are
  worth mentioning: if some factor asks for more than it holds, $|W_j| < c_j$, then
  $\binom{W_j}{c_j} = \emptyset$ and $\llbracket wl \rrbracket$ is empty; if every
  $c_j = 0$, the sole member is the empty probe. Either way there is nothing to
  certify. The verification task itself is the single-factor worklist
  $\big\langle (d, \mathcal{W}) \big\rangle$, whose denotation is
  $\binom{\mathcal{W}}{d}$, i.e. all $d$-subsets of the circuit's wires
  $\mathcal{W}$.

  Algorithm~\ref{alg:subset-space} sweeps this representation with three moves:
  discharge the union of the whole worklist and, on success, prune the branch;
  otherwise certify a representative probe and extend it to a safe set; then split
  each factor about that safe set and recurse on the pieces. The next two
  subsections treat the moves prune-and-extend and split, and presents a proof
  that the split loses nothing. In this section the extension part is kept ambigious,
  as its inner workings will be descredin in Section~\ref{sec:extension}.

\begin{algorithm}
\caption{Traversing the subset space (Component II)}
\label{alg:subset-space}
\begin{algorithmic}[1]
\Function{CheckAll}{$wl = \langle (c_1, W_1), \dots, (c_m, W_m) \rangle$}
  \If{$\llbracket wl \rrbracket = \emptyset$} \Return \Secure{} \Comment{some $|W_j| < c_j$, or every $c_j = 0$} \EndIf
  \State $U \gets W_1 \cup \dots \cup W_m$
  \State $(v, T) \gets \Call{CheckProbe}{U}$ \Comment{large-set pruning (\S\ref{sec:large-set-pruning})}
  \If{$v = \Secure{}$} \Return \Secure{} \EndIf
  \State $\mathbf{w} \gets$ a representative of $\llbracket wl \rrbracket$ \Comment{$c_j$ wires from each $W_j$}
  \State $(v, T) \gets \Call{CheckProbe}{\mathbf{w}}$ \Comment{Component I (\S\ref{sec:rewrite-engine})}
  \If{$v = \Insecure{}$} \Return \Insecure{}($\mathbf{w}$) \EndIf
  \State $\mathsf{safe} \gets \Call{Extend}{\mathbf{w}, T, U}$ \Comment{Component III (\S\ref{sec:extension}); \ $\mathbf{w} \subseteq \mathsf{safe} \subseteq U$}
  \ForAll{$(i_1, \dots, i_m) \in \textstyle\prod_{j} \{0, \dots, c_j\}$ \textbf{with} $(i_j)_j \neq (c_j)_j$}
    \State $wl' \gets \big\langle\, (i_j,\ W_j \cap \mathsf{safe}),\ \ (c_j - i_j,\ W_j \setminus \mathsf{safe}) \,\big\rangle_{j=1}^{m}$ \Comment{drop zero-count parts}
    \State $r \gets \Call{CheckAll}{wl'}$
    \If{$r \neq \Secure{}$} \Return $r$ \EndIf
  \EndFor
  \State \Return \Secure{}
\EndFunction
\end{algorithmic}
\end{algorithm}

  \subsection{Large-Set Pruning}
  \label{sec:large-set-pruning}

  Before splitting anything, the traversal tries to discharge the whole batch at
  once. Let $U = W_1 \cup \dots \cup W_m$ be every wire the worklist can reach,
  and hand $U$ --- as a single large probe --- to the rewrite engine of Component I. If
  $U$ is certified secret-independent, so is every subset of $U$, and in
  particular every probe in $\llbracket wl \rrbracket$; the branch is secure and
  the traversal returns without descending. On some secure circuits this one test
  does almost all of the work, some branches never split, because their union
  already discharges.

  However this discharge may fail, and it is worth seeing why the algorithm cannot stop here on
  failure. The union $U$ is typically far larger than $d$, and a large tuple can
  be secret-\emph{dependent} while every one of its size-$d$ sub-tuples is
  secret-independent, hence failing to discharge $U$ says nothing about the
  individual probes. When the union fails, we descend. We pick a
  representative $\mathbf{w} \in \llbracket wl \rrbracket$ --- one probe of
  the target size, $c_j$ wires from each factor --- certify it with Component I,
  and, if it holds, grow it with Component III (Section~\ref{sec:extension}) into a
  \emph{safe set}
  \[
    \mathsf{safe} \;\supseteq\; \mathbf{w}, \qquad
    \text{every subset of which is secret-independent.}
  \]
  This will be further explained in Component III. Put briefly,
  one measure-preserving coupling
  blinds every wire of $\mathsf{safe}$ simultaneously, so any subset of
  $\mathsf{safe}$, of any size, is secret-independent. Because the
  representative draws at least one wire from every factor, $\mathsf{safe}$ meets
  every $W_j$. This single fact is what makes the split below both exhaustive and
  strictly progressing.

  \subsection{Completeness via Vandermonde's Identity}
  \label{sec:vandermonde}

  With a safe set in hand the split is forced by a counting identity. Fix
  $\mathsf{safe}$ and cut each factor's ground set along it,
  \[
    W_j = A_j \uplus B_j, \qquad A_j = W_j \cap \mathsf{safe}, \quad B_j = W_j \setminus \mathsf{safe}.
  \]

  \begin{lemma}[Set Vandermonde]
  \label{lem:set-vandermonde}
  For disjoint wire sets $A, B$ and any $c \ge 0$,
  \[
    \binom{A \uplus B}{c} \;=\; \biguplus_{i=0}^{c}
      \Big\{\, S_A \cup S_B \;:\; S_A \in \tbinom{A}{i},\ S_B \in \tbinom{B}{\,c-i} \,\Big\},
  \]
  a \emph{disjoint} union: each $c$-subset $S$ lies in exactly the block
  $i = |S \cap A|$. Counting both sides recovers Vandermonde's identity,
  $\binom{|A|+|B|}{c} = \sum_{i=0}^{c} \binom{|A|}{i} \binom{|B|}{\,c-i}$.
  \end{lemma}

  The lemma is immediate, a $c$-subset is determined by its safe part and its
  unsafe part, chosen independently. Stating it set-wise, and not merely as the
  numeric identity, is the trick. It says not only how many probes each
  block holds but exactly which. Applying it in every factor and
  multiplying out,
  \[
    \llbracket wl \rrbracket \;=\; \prod_{j=1}^{m} \binom{W_j}{c_j}
    \;=\; \biguplus_{(i_1, \dots, i_m)} \ \llbracket wl_{(i_j)} \rrbracket,
    \qquad
    wl_{(i_j)} = \big\langle (i_j, A_j),\ (c_j - i_j, B_j) \big\rangle_{j=1}^{m},
  \]
  one \emph{cell} for each distribution $(i_1, \dots, i_m) \in \prod_j \{0, \dots, c_j\}$
  recording, factor by factor, how many of a probe's wires come from the safe
  side. The cells are disjoint and exhaust $\llbracket wl \rrbracket$, since every
  probe declares in each factor exactly how much of it is safe and each cell
  is again the denotation of a worklist, on which the traversal recurses directly.

  One cell needs no recursion however. The all-safe cell $(i_j)_j = (c_j)_j$
  collects the probes drawn entirely from the safe side $A_j \subseteq
  \mathsf{safe}$ and every one of
  those is secret-independent by the guarantee of Component III. The algorithm
  skips it and recurses on the rest. Skipping that particular cell means, with each
  recursion we are strictly decreasing the amount of probe sets we have to consider,
  which gives us our termination guarantee.

  \begin{theorem}[Soundness of the traversal]
  \label{thm:traversal}
  \textup{\textproc{CheckAll}}$(wl)$ (Algorithm~\ref{alg:subset-space})
  terminates, and if it returns \textup{\Secure{}} then every probe in
  $\llbracket wl \rrbracket$ is secret-independent. In particular
  \textup{\textproc{CheckAll}}$\big\langle (d, \mathcal{W}) \big\rangle$, returning
  \textup{\Secure{}}, certifies $d$-probing security.
  \end{theorem}

  \begin{proof}
  \emph{Termination.} Order worklists by the multiset $\{ |W_1|, \dots, |W_m| \}$
  of their ground-set sizes, under the multiset order. Take any cell the algorithm
  recurses on: it is neither the all-safe cell nor empty, so some factor
  $j^\ast$ has $i_{j^\ast} < c_{j^\ast}$, forcing $B_{j^\ast} \neq \emptyset$;
  and $A_{j^\ast} \supseteq \mathbf{w} \cap W_{j^\ast} \neq \emptyset$ because the
  representative meets every factor. So $W_{j^\ast}$ is replaced by the two
  strictly smaller ground sets $A_{j^\ast}, B_{j^\ast}$, while no other factor's
  ground set grows ($A_j, B_j \subseteq W_j$). The multiset strictly decreases, so
  the recursion is well-founded.

  \emph{Soundness}, by induction along that order. \textproc{CheckAll} returns
  \Secure{} only in one of three ways. If the union $U = \bigcup_j W_j$
  discharges, every subset of $U$ (hence all of $\llbracket wl \rrbracket$)
  is secret-independent by marginalisation. If a degenerate worklist bottoms out,
  $\llbracket wl \rrbracket$ is empty or the lone empty probe, which is trivially secure.
  Otherwise the representative certifies and the recursive cells all return
  \Secure{}. By Lemma~\ref{lem:set-vandermonde} those cells together with the
  skipped all-safe cell partition $\llbracket wl \rrbracket$; the all-safe cell is
  secret-independent by Component III, and each recursive cell is
  $\llbracket wl_{(i_j)} \rrbracket$ of strictly smaller measure, hence secret-independent
  by the induction hypothesis. Every probe of $\llbracket wl \rrbracket$ lies in
  exactly one part, so all are secret-independent. Finally, since
  $\llbracket \langle (d, \mathcal{W}) \rangle \rrbracket = \binom{\mathcal{W}}{d}$,
  a \Secure{} verdict on $\langle (d, \mathcal{W}) \rangle$ certifies that every $d$-subset
  of $\mathcal{W}$ is secret-independent, i.e.\ $d$-probing secure.
  \end{proof}

  We close this section with two remarks. The soundness above is one-directional by design:
  a \Secure{} verdict is trustworthy, but the traversal inherits Component
  I's incompleteness verbatim and nothing more. An \Insecure{} verdict is
  raised only when Component I declines a representative
  (Section~\ref{sec:general-substitution}), never because a subset slipped through
  unexamined. Lemma~\ref{lem:set-vandermonde} guarantees the cells cover
  $\llbracket wl \rrbracket$ exactly, so the sweep is genuinely exhaustive. The
  reported witness is therefore a true counterexample precisely when Component I is
  complete on it, as in the linear/affine case; a fuel-bounded refusal there is the
  one place a spurious \Insecure{} can enter, and it enters through Component
  I, not through this traversal. Second, the efficiency of the whole scheme rests
  on the safe set being large: the bigger $\mathsf{safe}$, the fewer and smaller
  the surviving cells, and the closer the union prune comes to settling a branch
  outright. That is exactly the quantity Component III's coupling extension is built
  to maximise, and it is where the approach of this thesis departs from a plain
  set-splitting search.

\section{Component III: Extending Probe Sets}
\label{sec:extension}

Component~II leans on a technique which we are yet to discuss. Every time it certifies
a representative probe $\mathbf{w}$, it asks for a whole \emph{safe set} of wires
around $\mathbf{w}$, every subset of which is again secret-independent, and it uses
that set to skip the all-safe cell of its Vandermonde split
(Section~\ref{sec:vandermonde}). Component~III is all about that safe set. It takes
the certified probe and grows it into such a safe set, reusing the work the rewrite
engine has already spent on $\mathbf{w}$ instead of certifying each new wire from
nothing. The larger the safe set it hands back, the more the traversal prunes, so
the whole scheme is only as efficient as this component is generous.

The three mechanisms we describe share one shape. Each of them walks over the wires
of the circuit that are not yet in the safe set and decides, for a candidate $u$,
whether the set stays jointly secret-independent once $u$ is added. What separates
them is the certificate they reuse and the test they run on it. Their soundness
rests on a single fact which is the following: suppose $\varphi$
is a measure-preserving bijection of the randomness space, and suppose that for
every wire $u$ in a set $U$ the composite $u \circ \varphi$ is independent of the secrets.
Since $\varphi$ preserves the uniform law, the joint distribution of the wires of
$U$ agrees with the joint distribution of the composites $u \circ \varphi$, and the
latter does not move with the secrets.
So the joint law of $U$ does not vary with the secrets, and by marginalisation neither
does the law of any subset of $U$. A single measure-preserving map that blinds
every wire of $U$ at once thus certifies the entire set. We should add that the
certified probe $\mathbf{w}$ is always kept in the safe set without a fresh test of course;
the per-wire tests below are run only on the other candidates. We present the three
mechanisms from the weakest to the strongest and adopt the last.

  \subsection{Closure-Driven Extension}
  \label{sec:closures}

  The first mechanism is set apart from the other two by \emph{when} it runs. A
  closure is not computed after a probe has been certified. It is grown together
  with the probe, in the same topological pass that assembles it. As the traversal
  walks the circuit and picks wires to fill the chosen probe of the current
  worklist, it keeps a running set of the wires already forced by the ones picked
  so far, and each new choice drains that set forward. Two local rules drive the
  drain. A gate all of whose inputs are already known becomes known in turn, since
  its value is then fixed, and when a known wire is an $\oplus$ of inputs of which
  all but one are known, that last input becomes known as well, recovered as the
  sum of the wire and the others. The pass finishes with the chosen probe and,
  around it, every wire these two rules could deduce from it.

  Everything the drain adds is a deterministic function of the probe. Writing
  $\mathbf{w}$ for the chosen probe and $u$ for a deduced wire, we have $u =
  f(\mathbf{w})$ for the function $f$ that the rules compose, so $u$ carries no
  randomness the probe does not already carry, i.e. $H(u \mid \mathbf{w}) = 0$.
  Adjoining it leaves the joint entropy untouched, $H(\mathbf{w} \cup \{u\}) =
  H(\mathbf{w})$, and the pair $(\mathbf{w}, u)$ is a deterministic image of
  $\mathbf{w}$. Its dependence on the secrets is therefore exactly the probe's,
  which the engine has already certified to be none. That is the whole of the
  closure argument, and it needs no coupling.

  What the closure buys is little. In a masked circuit almost no wire is a plain
  function of a handful of others, so the set deduced around a small probe is
  rarely more than a few relabellings, and the traversal is left to split nearly
  the whole space by hand. The shortfall sharpens as Component~II recurses. Its
  worklist factors draw from wire sets carved out of the safe and unsafe splits, so
  deep in the recursion the candidate pool is a scattering of wires gathered from
  distant, unrelated parts of the circuit. Closure deduces a wire only from its
  immediate neighbours, and those neighbours become rare in a fragmented pool,
  so the deduction finds little to do exactly where the traversal needs
  the help most. In our experiments this mechanism extended probes the least of the
  three, although it is the cheapest of them, since the deduction rides along inside
  a pass the traversal was making anyway and never calls the engine a second time.
  Its worth is mostly conceptual.

  \subsection{Replay-Driven Extension}
  \label{sec:replay}

  A stronger idea, due to Barthe et al.\ \parencite{10.1007/978-3-030-29959-0_15},
  is to reuse the derivation rather than the probe. Certifying $\mathbf{w}$ left a
  short record of the steps that discharged it (Section~\ref{sec:coupling-view}),
  and each step is an operation one can carry out on any expression, not only on
  $\mathbf{w}$. Replay-driven extension takes a candidate wire $u$ from the pool of
  the current worklist and runs that same record on it, step by step. If $u$ comes
  out free of secrets at the end, then it is blinded by the very derivation that
  blinded the probe, and it joins the safe set. This reaches well past the closure,
  because a wire no longer has to be a function of the probe to survive the replay.
  It only has to be carried to a secret-free form by the same steps, which many of
  the sibling wires of a masked gadget are.

  Barthe et al.'s derivations are built from two kinds of step.
  The first is the substitution of an invertible random by its context, the
  optimistic-sampling move our own engine makes. The second is an algebraic
  normalisation, which rewrites an expression into its algebraic normal form (its
  unique multilinear $\mathbb{F}_2$ polynomial, Section~\ref{sec:field-f2}) and so
  lets products cancel. Barthe turns to it only when no substitution applies, and
  records it in the derivation like any other step. When the replay meets such a
  step it normalises the candidate too.

  For us replays are a little problematic. Our engine does not lean on
  normalisation to free a buried random. It factors the expression instead
  (Section~\ref{sec:factoring}), reshaping a product so that a random trapped
  inside it moves into an additive position where a substitution can reach it.
  Factoring wins the engine circuits that substitution and normalisation alone
  would miss, but it is not an easily replayble step. Re-applying a factoring
  to an unrelated candidate wire does not carry over the blinding the way
  re-applying a substitution does, so the record we keep holds only the
  substitutions and leaves the factoring implicit. Replaying that record on a
  candidate is then incomplete. A wire the derivation genuinely blinds can fail it,
  because the cancellation that clears its secret was arranged by a factoring the
  replay never performs. That mismatch is the whole
  reason why we take one further step, and read the derivation not as a script to rerun
  but as a coupling function.

  \subsection{Coupling-Driven Extension}
  \label{sec:coupling-extension}

  The mechanism we adopt keeps replay's use of the derivation and replaces its
  final test with an exact one. Section~\ref{sec:coupling-view} already read the
  recorded derivation as a single measure-preserving bijection $\varphi$ of the
  randomness space, the \gls{coupling} that certified $\mathbf{w}$. Coupling-driven
  extension admits a candidate wire $u$ exactly when $u \circ \varphi$ is
  independent of the secrets, and it decides that independence semantically rather
  than by surface syntax. Every wire replay
  admits is admitted here too, and so are the wires replay had to withhold because a
  factoring had to set up cancellations.

  The candidates are not all the wires of the
  circuit. They are the wires of the current worklist, the pool $C$ that
  Component~II is trying to certify (Section~\ref{sec:subset-space}), and the probe
  $\mathbf{w}$ is kept among them with no test at all. To decide $u \circ
  \varphi$ is secret-free exactly, for every candidate, is not cheap, so the work is
  staged. Initially, a fast probabilistic filter throws out the candidate wires
  that are incompatible with the coupling function we derived via evaluations.

\begin{algorithm}
\caption{Coupling-driven extension (Component III)}
\label{alg:extend}
\begin{algorithmic}[1]
\Function{Extend}{$\mathbf{w}, \varphi, C$}
  \Require certified probe $\mathbf{w}$, coupling $\varphi$, candidate pool $C$ from the worklist
  \Ensure safe set $\mathsf{safe}$ with $\mathbf{w} \subseteq \mathsf{safe} \subseteq C$, with every subset secret-independent
  \State $\mathsf{safe} \gets \mathbf{w}$
  \ForAll{$u \in C \setminus \mathbf{w}$}
    \If{\Call{FastReject}{$u, \varphi$}}
      \State \textbf{continue} \Comment{Phase 1: probabilistic filter, rejects proven-dependent $u$ only}
    \EndIf
    \If{\Call{SurfaceFree}{$u, \varphi$} \textbf{ or } \Call{AnfFree}{$u, \varphi$}} \Comment{Phases 2--3: exact, decides $u \circ \varphi$ secret-free}
      \State $\mathsf{safe} \gets \mathsf{safe} \cup \{u\}$
    \EndIf
  \EndFor
  \State \Return $\mathsf{safe}$
\EndFunction
\end{algorithmic}
\end{algorithm}

  The filter is an evaluation under two runs. Fix random values for the
  randoms, evaluate $u \circ \varphi$ once with the secrets also drawn at random and
  once with the secrets set to zero, and compare. If the two disagree, then $u \circ
  \varphi$ genuinely moves with the secrets, and $u$ is dropped. Agreement proves
  nothing on its own, since a dependence can hide from any one tape, so a survivor
  is not yet admitted. Running the trial over many tapes makes a surviving
  dependence steadily less likely but never impossible, which is why the filter is
  allowed to reject but never to accept.

  Then the derivation is replayed on $u$, and the
  result is first read by ordinary $\mathbb{F}_2$ cancellation, which is cheap and
  clears the common case. When a secret still shows, the wire is converted to its
  algebraic normal form and admitted exactly when no secret variable survives there.
  The probabilistic filter and this exact stage differ in kind, so it is worth
  recording that together they never admit a leaking wire and never turn away a
  blinded one.

  \begin{lemma}[The extension test is exact]
  \label{lem:extend-exact}
  \textproc{Extend} admits a candidate $u$ if and only if $u \circ \varphi$ is
  independent of the secrets.
  \end{lemma}

  \begin{proof}
  Write $v = u \circ \varphi$. Its algebraic normal form (ANF) is the unique
  multilinear $\mathbb{F}_2$ polynomial computing $v$, so a variable is absent from
  that form exactly when $v$ does not depend on it. In particular,
  \[
    v \text{ is independent of the secrets}
    \iff
    \text{no secret variable occurs in the ANF of } v .
  \]
  A wire is admitted only at the exact stage. Phase~2 accepts when the surface form
  of $v$ carries no secret, and Phase~3 accepts when the ANF of $v$ carries no
  secret variable. By the equivalence, each makes $v$ independent of the secrets, so
  every admitted wire is truly blinded. Conversely, suppose $v$ is independent of the
  secrets. With the randoms fixed, $v$ is then constant in the secrets, so the
  filter's two evaluations agree on every tape and never reject $u$, and it reaches
  the exact stage. And because the ANF of $v$ must not carry a secret variable, Phase~3
  will admit it. Hence $u$ is admitted if and only if $v$ is independent of the secrets.
  \end{proof}

  \begin{theorem}[Safe set]
  \label{thm:safe-set}
  Let $\mathsf{safe} = \textproc{Extend}(\mathbf{w}, \varphi, C)$ for a probe
  $\mathbf{w}$ certified by Component~I with coupling $\varphi$. Then the joint
  distribution of the wires of $\mathsf{safe}$, over a uniformly random $\rho$, does
  not depend on the secret assignment, and in particular every subset of
  $\mathsf{safe}$ is secret-independent.
  \end{theorem}

  \begin{proof}
  By Lemma~\ref{lem:extend-exact}, $\mathsf{safe}$ consists of the probe $\mathbf{w}$
  and the candidates $u$ whose composite $u \circ \varphi$ is independent of the
  secrets. Since that composite does not depend on the secret assignment, it is a
  function of the randomness alone; so simply write its value as $\overline{u}(\rho)$.

  Fix a secret assignment $\mathbf{a}$ and let $\varphi_{\mathbf{a}}$ be the induced
  measure-preserving bijection of the randomness space
  (Section~\ref{sec:coupling-view}). Chaining Lemma~\ref{lem:substitution} through
  the recorded steps gives, for every $u \in \mathsf{safe}$ and every $\rho$,
  \[
    u\big(\varphi_{\mathbf{a}}(\rho);\, \mathbf{a}\big)
    \;=\; (u \circ \varphi)(\rho;\, \mathbf{a})
    \;=\; \overline{u}(\rho) .
  \]
  The left-hand side is $u$ evaluated at the relabelled randomness assignment
  $\varphi_{\mathbf{a}}(\rho)$ under the secrets $\mathbf{a}$, so it is a bit value;
  the first equality is the substitution lemma, the second one comes from the
  admission condition. Both copies of
  $\mathbf{a}$, the one inside $\varphi_{\mathbf{a}}$ and the one feeding $u$ its
  secret inputs, cancel, and the value depends on $\rho$ alone.

  Now hold $\mathbf{a}$ fixed and draw $\rho$ uniformly. As $\varphi_{\mathbf{a}}$
  preserves the uniform law, $\varphi_{\mathbf{a}}(\rho)$ is uniform too, so the
  joint law of the safe wires under $\mathbf{a}$, that of
  $\big(u(\rho;\, \mathbf{a})\big)_{u \in \mathsf{safe}}$, equals the joint law of
  $\big(u(\varphi_{\mathbf{a}}(\rho);\, \mathbf{a})\big)_{u} =
  \big(\overline{u}(\rho)\big)_{u}$. This last family depends on no secret, so the
  joint law is one and the same for every $\mathbf{a}$. Hence the distribution of
  $\mathsf{safe}$ does not depend on the secret assignment, nor any of its subsets.
  \end{proof}

  Finally we observe that the other two mechanisms admit nothing if this one does not.

  \begin{proposition}[Closure and replay sit inside the coupling extension]
  \label{prop:nesting}
  Fix a certified probe $\mathbf{w}$ with coupling $\varphi$, and write
  $\mathsf{safe}_{\mathrm{clo}}$, $\mathsf{safe}_{\mathrm{rep}}$, and
  $\mathsf{safe}_{\mathrm{cou}}$ for the wires admitted by the closure, replay, and
  coupling extensions. Then $\mathsf{safe}_{\mathrm{clo}} \subseteq
  \mathsf{safe}_{\mathrm{cou}}$ and $\mathsf{safe}_{\mathrm{rep}} \subseteq
  \mathsf{safe}_{\mathrm{cou}}$.
  \end{proposition}

  \begin{proof}
  For replay, a wire is admitted only after its replay rechecks to a secret-free
  form, which means $u \circ \varphi$ is independent of the secrets; the coupling
  test admits every such wire by Lemma~\ref{lem:extend-exact}. For the closure, an
  admitted wire is a gate-composition $u = f(\mathbf{w})$. Evaluation respects that
  composition, so $u \circ \varphi = f(\mathbf{w} \circ \varphi)$ as functions of the
  randomness; the coupling leaves $\mathbf{w} \circ \varphi$ independent of the
  secrets, and a gate-composition of functions independent of the secrets is again
  independent of them, so $u \circ \varphi$ is too and the coupling test admits $u$.
  \end{proof}
